\section{Discussion}

This study developed a machine learning-based diabetes prediction system using the Philippine Expanded National Nutrition Survey (ENNS) 2013 data. The system integrates multiple machine learning algorithms, probability calibration techniques, and explainable artificial intelligence (XAI) methods to provide interpretable and reliable diabetes risk predictions. The use of XGBoost combined with SHAP for feature importance analysis and Platt Scaling for probability calibration provided both high predictive performance and clinical interpretability.

\subsection{Model Selection and Performance}

While CatBoost and LightGBM achieved higher precision (0.558 and 0.552) and ROC-AUC (0.716 and 0.710) scores, they demonstrated significantly lower recall (0.442 and 0.435) compared to XGBoost. This indicates that these models may miss more than half of diabetic cases despite their accuracy in predicting true positives. In a diabetes screening context, missing diabetic cases (false negatives) has more serious clinical consequences than false alarms (false positives), as undetected diabetes can lead to severe complications including cardiovascular disease, kidney failure, and neuropathy.

k-Nearest Neighbors (kNN) demonstrated the second-highest recall of 0.597, but its lower precision score of 0.388 indicated a higher rate of false positives. This would result in unnecessary confirmatory testing and patient anxiety. Therefore, XGBoost was selected as the optimal predictive model for diabetes screening. XGBoost achieved the highest F2-score of 0.692 among all models and demonstrated consistent high scores across all performance metrics, balancing the need for high recall (80.1\%) with acceptable precision (44.9\%).

The ensemble methods (Voting and Stacking) did not outperform XGBoost despite combining multiple base models. The Voting Ensemble achieved an F2-score of 0.565, while the Stacking Ensemble achieved 0.452, both substantially lower than XGBoost's 0.692. This suggested that XGBoost's performance was so dominant that averaging it with weaker models degraded overall performance rather than improving it. This finding contrasts with typical ensemble learning expectations, where combining diverse models usually improves performance. The dominance of XGBoost indicates that tree-based gradient boosting is particularly well-suited for tabular health data with mixed feature types.

TabNet, despite being a deep learning architecture specifically designed for tabular data with attention mechanisms, underperformed compared to XGBoost. The Bayesian Optimized TabNet achieved an F2-score of 0.570, which was 17.6\% lower than XGBoost. The extensive hyperparameter optimization process (12 iterations $\times$ 5 folds = 60 evaluations, 214 minutes) provided minimal improvement over the baseline TabNet (F2-score increase of only 0.007). This suggests that deep learning models may not be necessary or beneficial for datasets of this size (approximately 67,000 training samples). Deep learning typically requires hundreds of thousands to millions of samples to fully leverage its capacity, and the relatively modest dataset size may have limited TabNet's ability to learn complex patterns. Additionally, the attention mechanism in TabNet may have focused on spurious correlations rather than genuine clinical risk factors, reducing its generalization capability.

\subsection{Feature Importance and Clinical Validity}

The SHAP analysis revealed that the most influential features for diabetes prediction were average systolic blood pressure (Ave\_SBP), waist circumference, and triglycerides. These top three features are well-established risk factors for Type 2 diabetes in medical literature and align with the components of metabolic syndrome. Metabolic syndrome is a cluster of conditions including hypertension, central obesity, dyslipidemia, and insulin resistance that collectively increase the risk of Type 2 diabetes and cardiovascular disease.

Blood pressure, particularly systolic blood pressure, is strongly associated with insulin resistance and diabetes risk. The dominance of Ave\_SBP as the top predictor (SHAP value of 0.0564) confirms its critical role in diabetes pathophysiology. Waist circumference, the second most important feature (SHAP value of 0.0499), is a measure of central obesity and visceral fat accumulation. Visceral fat is metabolically active and contributes to insulin resistance through the release of inflammatory cytokines and free fatty acids. Triglycerides, the third most important feature (SHAP value of 0.0098), indicate lipid metabolism disorders commonly associated with insulin resistance and prediabetes.

Anthropometric features (waist, weight, hip) appeared prominently in the top 10 features, highlighting the importance of body composition in diabetes prediction. Dietary features (food groups 15 and 25, energy-weighted intakes) also contributed to predictions, though with lower SHAP values than clinical and anthropometric measurements. This suggests that while dietary patterns play a role in diabetes risk, their predictive power is lower than direct physiological measurements.

Notably, all top features represent clinically valid health indicators rather than administrative codes or survey design variables. This contrasts with models that may learn spurious patterns from geographic codes (region, province) or sampling design variables. The removal of administrative variables (\texttt{regcode}, \texttt{provcode}, \texttt{psurec}, \texttt{strrec}) ensured that the model learned genuine clinical risk factors rather than geographic or demographic proxies. This makes the model more generalizable across different Philippine regions and populations.

The consistency between SHAP importance and XGBoost's native feature importance further validated the robustness of the results. Both methods identified Ave\_SBP and waist as the top two most important features, with only minor differences in ranking for other features. This agreement indicates that the model's decision-making process is stable and interpretable, which is critical for clinical adoption and trust.

\subsection{Probability Calibration and Clinical Utility}

Probability calibration using Platt Scaling improved the Brier Score from 0.2364 to 0.2016, representing a 14.7\% improvement. This calibration ensures that predicted probabilities accurately reflect true diabetes risk. For example, when the calibrated model predicts a 70\% probability of diabetes, approximately 70\% of such cases are actually diabetic. Without calibration, the predicted probabilities may be systematically overconfident or underconfident, leading to unreliable risk assessments.

The calibration improvement is particularly important for clinical decision-making and risk-based screening strategies. Calibrated probabilities enable healthcare providers to stratify patients into risk categories (low, moderate, high) and allocate resources accordingly. High-risk patients can be prioritized for confirmatory testing and early intervention, while low-risk patients can be monitored through regular screening. This risk-based approach is more efficient than universal testing and allows for personalized diabetes prevention strategies.

The threshold analysis demonstrated that the model's performance could be tuned based on the specific clinical application. At a threshold of 0.3, the model achieved 87.6\% recall, making it suitable for maximum sensitivity screening where the goal is to catch as many diabetic cases as possible. At a threshold of 0.7, the model achieved 53.9\% precision, making it more appropriate for confirmatory testing where reducing false positives is prioritized. This flexibility allows the model to be adapted to different healthcare settings and resource constraints.

\subsection{Integration of Explainable Artificial Intelligence}

The integration of explainable AI techniques such as SHAP enhanced the transparency of the model's decision-making process. This addressed the criticized "black-box" nature of machine learning models, which has been a major barrier to clinical adoption. By providing feature-level explanations showing which factors contribute to each prediction, SHAP helps foster trust and understanding among healthcare professionals.

SHAP provides both global feature importance (ranking features by their overall impact across all predictions) and local explanations (showing how features contribute to individual predictions). Global feature importance helps identify population-level risk factors and validate that the model is learning clinically meaningful patterns. Local explanations enable personalized risk communication, allowing healthcare providers to explain to individual patients which specific factors (e.g., high blood pressure, large waist circumference) are driving their diabetes risk.

The use of calibrated probabilities in conjunction with SHAP explanations provides a comprehensive framework for clinical decision support. Healthcare providers can see not only the predicted risk level but also the specific factors contributing to that risk, along with uncertainty estimates. This combination of prediction, explanation, and uncertainty quantification addresses the key requirements for trustworthy AI in healthcare.

\subsection{Comparison with Reference Study}

This study differs from the reference work by Rivera (2024) in several important methodological aspects. First, this study merged all four ENNS component datasets (clinical, biochemical, anthropometric, dietary), resulting in a more comprehensive feature set of 64 features compared to Rivera's 47 features. The inclusion of dietary data provides additional information about nutritional patterns and their relationship with diabetes risk.

Second, this study applied SMOTE for class imbalance handling \textbf{after} the train-validation-test split to prevent data leakage. Applying SMOTE before splitting would allow synthetic instances based on test data to appear in the training set, artificially inflating performance metrics. Rivera's methodology applied SMOTE before splitting, which may have contributed to optimistic performance estimates. The proper application of SMOTE in this study ensures that the reported metrics reflect true generalization performance.

Third, this study removed all administrative and survey design variables (\texttt{regcode}, \texttt{provcode}, \texttt{psurec}, \texttt{strrec}, sampling weights) to prevent data leakage. These variables are artifacts of the survey design and do not represent individual health characteristics. Models that include these variables may learn geographic or demographic patterns rather than clinical risk factors, reducing generalizability to other populations. The removal of these variables ensures that the model relies solely on clinically valid health indicators.

Fourth, this study tested a broader range of models (13 different approaches including traditional ML, deep learning, and ensembles) compared to Rivera's focus on TabNet and five traditional ML models. This comprehensive comparison provides stronger evidence for the superiority of XGBoost and demonstrates that neither deep learning nor ensemble methods provide advantages for this particular dataset and task.

The F2-score achieved in this study (0.692) is lower than Rivera's reported F2-score (0.842). However, this difference is attributable to the methodological improvements implemented in this study, particularly the proper application of SMOTE after splitting and the removal of administrative variables. The lower but more rigorous metrics in this study represent true generalization performance rather than inflated estimates from data leakage. The 80.1\% recall achieved in this study is still appropriate for diabetes screening applications, where the goal is to catch the majority of diabetic cases while maintaining acceptable precision.

\subsection{Challenges and Limitations}

Several challenges were encountered during the development of this diabetes prediction system. The class imbalance in the dataset required careful handling through SMOTE, but even with oversampling, the minority class remained underrepresented. This contributed to the moderate precision (44.9\%) and highlights the inherent difficulty of predicting diabetes in population-level surveys where the prevalence is relatively low.

The missing values in the ENNS dataset, particularly in dietary features, required extensive imputation. While KNN imputation preserves local structure better than simple mean imputation, it still introduces uncertainty into the data. Features with more than 50\% missing values were dropped entirely to avoid excessive imputation bias, but this resulted in the loss of potentially informative variables.

The cross-sectional nature of the ENNS data limits the ability to establish causal relationships between risk factors and diabetes. The associations identified through SHAP analysis are correlational rather than causal, and temporal relationships cannot be determined. Longitudinal data would be needed to track how risk factors evolve over time and predict diabetes onset.

The dataset is limited to a single survey period (2013), which may not reflect current dietary patterns, lifestyle behaviors, or healthcare access in the Philippines. External validation on more recent ENNS data or independent Philippine datasets would be needed to assess the model's generalizability across time periods.

\subsection{Main Contributions}

This study makes several important contributions to diabetes prediction research in the Philippine context. First, it provides the first rigorous integration of all four ENNS component datasets (clinical, biochemical, anthropometric, dietary) for diabetes prediction, resulting in a more comprehensive feature set than previous studies. This multi-domain approach captures the complex interplay between physiological, anthropometric, and dietary factors in diabetes risk.

Second, the study identifies and corrects methodological issues in existing approaches, particularly data leakage from improper SMOTE application and inclusion of administrative variables. The documentation of these issues and their impact on performance metrics (5-7\% inflation) provides valuable guidance for future researchers working with survey data.

Third, the comprehensive comparison of 13 different modeling approaches demonstrates that tree-based gradient boosting (XGBoost) outperforms both deep learning (TabNet) and ensemble methods for Philippine health data. This finding challenges the assumption that more complex models always perform better and highlights the importance of matching model complexity to dataset characteristics.

Fourth, the integration of probability calibration with explainable AI provides a framework for trustworthy clinical decision support. The combination of calibrated risk predictions, SHAP-based feature explanations, and uncertainty quantification addresses the key requirements for AI adoption in healthcare settings.

Finally, the study demonstrates that high recall (80.1\%) for diabetes screening can be achieved using readily available health measurements (blood pressure, waist circumference, triglycerides) without relying on expensive laboratory tests or genetic information. This makes the approach feasible for resource-constrained healthcare settings in the Philippines and other developing countries.